\documentclass[10pt]{article}
\usepackage[T1]{fontenc}
\usepackage{lmodern,amsmath,amssymb,amsthm,microtype}
\usepackage[margin=0.82in]{geometry}
\usepackage[colorlinks=true,urlcolor=blue]{hyperref}
\newtheorem{theorem}{Theorem}
\newtheorem{lemma}{Lemma}
\newtheorem{corollary}{Corollary}
\title{Six Forbidden Prime Classes in an Exponential Sequence}
\author{Qichao Wang\\{\small Hebei University of Technology}\\[0.3em]
Daoyu Dong\\{\small University of Electronic Science and Technology of China}}
\date{October 26, 2026}
\begin{document}
\maketitle
\begin{abstract}
We prove the six excluded prime residue classes conjectured in
OEIS A237930. Each excluded prime makes three a quadratic residue
and seven a quadratic nonresidue, which prevents it from dividing
any term $(7\cdot3^n-1)/2$. We also give the exact finite
multiplicative-order criterion for every other odd prime.
\end{abstract}

\section{The six forbidden classes}
Write
\[
 A(n)=3^{n+1}+\frac{3^n-1}{2}=\frac{7\cdot3^n-1}{2},\qquad n\ge0.
\]
The April 13, 2020 comment in OEIS A237930 conjectures that no
prime congruent to $11,13,23,61,71,$ or $73$ modulo $84$ divides
any term. The listed classes are interpreted as six integers;
the missing comma between $61$ and $71$ in the entry is typographical.

For an odd prime $p$ and integer $a$ coprime to $p$, denote the
Legendre symbol by $\left(\frac ap\right)$.
\begin{theorem}
For every $n\ge0$, no prime divisor of $A(n)$ is congruent to any of
\[
 11,13,23,61,71,73\pmod{84}.
\]
\end{theorem}
\begin{proof}
Quadratic reciprocity and its supplementary laws give
\[
 \left(\frac3p\right)=
 \begin{cases}1,&p\equiv1,11\pmod{12},\\-1,&p\equiv5,7\pmod{12},\end{cases}
 \qquad
 \left(\frac7p\right)=(-1)^{(p-1)/2}\left(\frac p7\right).
\]
The nonzero squares modulo seven are $1,2,4$. Applying these
rules to the unit classes modulo $84$ gives
\[
 \left(\frac3p\right)=1,\quad\left(\frac7p\right)=-1
 \quad\Longleftrightarrow\quad
 p\equiv11,13,23,61,71,73\pmod{84}.
\]
For transparency, the associated local residues are
\[
\begin{array}{c|rrrrrr}
 p\bmod84&11&13&23&61&71&73\\\hline
 p\bmod12&11&1&11&1&11&1\\
 p\bmod7&4&6&2&5&1&3\\
 p\bmod4&3&1&3&1&3&1
\end{array}
\]
If such a $p$ divided $A(n)$, then $7\cdot3^n\equiv1\pmod p$.
Taking Legendre symbols would give
\[
 1=\left(\frac7p\right)\left(\frac3p\right)^n=-1,
\]
which is impossible.
\end{proof}

\section{The exact incidence condition for other primes}
The obstruction proves precisely the six-class conjecture; it
need not characterize all primes that never divide the sequence.
An exact finite test for each remaining prime is as follows.
The prime $2$ divides $A(n)$ exactly when $n$ is odd, since
$7\cdot3^n-1\equiv0\pmod4$ exactly then. The prime $3$ divides
only $A(0)=3$, and $7$ divides no term.

For an odd prime $p\ne3,7$, put $t=\operatorname{ord}_p(3)$.
If the finite set
\[
 H_p=\{0\le h<t:3^h\equiv7^{-1}\pmod p\}
\]
contains an element $h$, it is unique, and
\[
 p\mid A(n)\quad\Longleftrightarrow\quad n\equiv h\pmod t.
\]
If $H_p$ is empty, $p$ divides no term. This follows immediately
from the periodic and distinct powers $3^0,\ldots,3^{t-1}$.
Thus the sequence's individual prime-divisibility questions are
also completely decidable by finite arithmetic.

\section{Independent checks}
The accompanying verifier first classifies all unit residues modulo
$84$, then checks the two quadratic symbols by the independent
Euler criterion for all relevant primes through $200000$.
This includes 4513 primes in the excluded classes.
For every excluded prime below $5000$, it independently traverses
the entire orbit of powers of three and checks absence of a
solution; these cycles contain 134995 residues in total.
All 1001 official sequence terms agree with the exponential formula
and the independent affine recurrence $A(0)=3$,
$A(n+1)=3A(n)+1$. The all-index exclusion is proved above,
not inferred from those finite tests.

\begin{thebibliography}{1}
\bibitem{oeis} The OEIS Foundation Inc.,
\href{https://oeis.org/A237930}{A237930} and its
\href{https://oeis.org/A237930/b237930.txt}{b-file},
accessed October 26, 2026.
\end{thebibliography}
\end{document}
